\chapter{Proposta}
\label{cap:proposta}

Este capítulo define o que foi proposto antes de apresentar como a implementação foi conduzida ou quais resultados foram obtidos. Trata-se de um protótipo acadêmico para exercitar o caminho entre dados sintéticos, software ARM e lógica FPGA. A proposta não inclui conexão a bolsa nem emissão de ordens.

\section{Arquitetura e fronteiras}\label{sec:propostaVisao}

A arquitetura separa o sistema em três domínios complementares. No software, um gerador publica mensagens FAST sintéticas por TCP e um receptor C++ as decodifica; quando o acesso ao FPGA está habilitado, esse receptor converte símbolo, lado, preço e quantidade para palavras inteiras. Na lógica programável, duas filas transportam eventos e respostas, um livro de ofertas mantém níveis por símbolo e uma estratégia produz um código de decisão. No fluxo de desenvolvimento, o HDL Coder gera o VHDL da estratégia a partir de uma função MATLAB.

A \autoref{fig:placeholderArquiteturaGeral} apresenta o percurso de execução e, pela ligação tracejada, o percurso de geração. Durante a execução, os eventos seguem do fluxo sintético para o ARM, atravessam a ponte MMIO e chegam ao FPGA; a resposta retorna ao receptor. MATLAB e HDL Coder não participam desse percurso em tempo de execução: eles produzem um arquivo VHDL que é incorporado ao projeto antes da síntese.

\begin{figure}[htb]
\caption{Arquitetura geral do protótipo}
\label{fig:placeholderArquiteturaGeral}
\centering
\includegraphics[height=0.58\textheight]{arquitetura-geral}
\fonte{Elaborado pelos autores}
\end{figure}

Esse percurso define a responsabilidade de cada bloco: o ARM trata protocolo e conversão de representação; a ponte e as filas implementam o contrato binário; o livro mantém o estado simplificado; e a estratégia devolve NOOP, BUY ou SELL. Esses códigos são saídas internas de teste, não ordens transmitidas a um mercado. A arquitetura cria três fronteiras verificáveis: bytes FAST tornam-se campos C++ tipados; esses campos tornam-se um quadro de oito palavras de 32 bits; e a atualização do livro produz os sinais de topo consumidos pela estratégia.

\section{Decisões arquiteturais}\label{sec:decisoesArquiteturais}

A decodificação FAST permanece no ARM porque envolve biblioteca, modelo de mensagem, campos textuais e tratamento de conexão. Com isso, o FPGA recebe uma representação já normalizada. A escolha reduz o escopo do hardware, mas mantém o sistema operacional no caminho de execução e afasta o protótipo de arquiteturas em que a FPGA recebe quadros diretamente da rede.

A comunicação usa duas filas circulares em MMIO: TX, do ARM para o FPGA, e RX, no sentido inverso. As filas desacoplam produtor e consumidor enquanto houver posições livres. A capacidade útil é uma posição menor que a profundidade física, e a publicação dos ponteiros só ocorre depois da escrita ou leitura completa do quadro.

Livro e estratégia são tratados como blocos distintos. O livro mantém estado e ordena até oito níveis de compra e oito de venda por símbolo; a estratégia, por sua vez, é combinacional e recebe melhor compra, melhor venda, respectivas quantidades, \textit{spread} e desbalanceamento. Uma alteração da regra exige nova geração, novos testes e nova síntese do projeto, mas não exige mudar o formato do evento enquanto a interface do bloco for preservada.

A \autoref{fig:placeholderLivro} mostra essa separação. O quadro é decodificado e direcionado ao símbolo e ao lado correspondentes; depois da atualização, os sinais do topo alimentam a estratégia e integram a resposta. A figura representa o fluxo lógico, não a temporização em ciclos.

\begin{figure}[htb]
\caption{Fluxo lógico do livro de ofertas simplificado}
\label{fig:placeholderLivro}
\centering
\includegraphics[height=0.66\textheight]{livro-ofertas}
\fonte{Elaborado pelos autores}
\end{figure}

O ponto substituível da proposta é a função \codigo{strategy.m}, com seis entradas inteiras e uma saída de decisão. A substituição preserva a interface com o livro, mas exige gerar novamente o VHDL, executar os testes e recompilar o projeto FPGA. O procedimento correspondente é descrito no \autoref{sec:metodo}.

\section{Contrato ARM--FPGA}\label{sec:contratoEventos}

Cada posição das filas contém oito palavras de 32 bits. A \autoref{tab:quadroEvento} apresenta o quadro de entrada. Os campos 6 e 7 são reservados e valem zero na implementação atual. O evento de reinicialização limpa o livro do símbolo indicado; o receptor do fluxo sintético publica eventos do tipo inserção ou atualização.

\begin{table}[htb]
\caption{Quadro enviado do ARM para o FPGA}
\label{tab:quadroEvento}
\begin{tabularx}{\textwidth}{lllX}
\toprule
\textbf{Palavra} & \textbf{Campo} & \textbf{Representação} & \textbf{Semântica} \\ \midrule
0 & \codigo{seq_no} & sem sinal, 32 bits & Número de sequência sintético. \\
1 & \codigo{symbol_id} & sem sinal, 32 bits & Identificador numérico do símbolo. \\
2 & \codigo{price_1e4} & sem sinal, 32 bits & Preço multiplicado por $10^4$. \\
3 & \codigo{qty} & sem sinal, 32 bits & Quantidade agregada no nível. \\
4 & \codigo{event_type} & sem sinal, 32 bits & 1: inserir/atualizar; 2: remover; 3: reinicializar. \\
5 & \codigo{side} & sem sinal, 32 bits & 1: compra; 2: venda. \\
6 & reservado & 32 bits & Reservado para extensão. \\
7 & reservado & 32 bits & Reservado para extensão. \\ \bottomrule
\end{tabularx}
\fonte{Elaborado pelos autores conforme o contrato implementado em C++ e VHDL}
\end{table}

A \autoref{tab:quadroResposta} descreve a resposta. O campo \codigo{action} é um código interno de decisão. BUY e SELL indicam qual ramo da regra foi ativado; nenhum dos dois provoca transmissão de ordem.

\begin{table}[htb]
\caption{Quadro devolvido do FPGA para o ARM}
\label{tab:quadroResposta}
\begin{tabularx}{\textwidth}{lllX}
\toprule
\textbf{Palavra} & \textbf{Campo} & \textbf{Representação} & \textbf{Semântica} \\ \midrule
0 & \codigo{seq_no} & sem sinal, 32 bits & Sequência ecoada para correlação. \\
1 & \codigo{action} & sem sinal, 32 bits & 0: NOOP; 1: BUY; 2: SELL. \\
2 & \codigo{best_bid_px} & sem sinal, 32 bits & Melhor preço de compra, em escala $10^4$. \\
3 & \codigo{best_bid_qty} & sem sinal, 32 bits & Quantidade da melhor compra. \\
4 & \codigo{best_ask_px} & sem sinal, 32 bits & Melhor preço de venda, em escala $10^4$. \\
5 & \codigo{best_ask_qty} & sem sinal, 32 bits & Quantidade da melhor venda. \\
6 & \codigo{spread_1e4} & sem sinal, 32 bits & Melhor venda menos melhor compra. \\
7 & \codigo{imbalance} & com sinal, 32 bits & Quantidade compradora menos vendedora. \\ \bottomrule
\end{tabularx}
\fonte{Elaborado pelos autores conforme o contrato implementado em C++ e VHDL}
\end{table}

O FPGA informa assinatura, versão, profundidades e número de palavras por posição, e o software usa essa geometria para calcular a base da fila RX por

\begin{equation}
  \mathrm{RX\_BASE} = \mathrm{TX\_BASE} + \mathrm{TX\_DEPTH}\times\mathrm{SLOT\_WORDS}\times 4.
\end{equation}

Com profundidade 64, oito palavras e base TX \codigo{0x100}, a base RX é \codigo{0x900}. O cálculo evita vincular o software a uma única geometria.

\section{Símbolos, regra de decisão e delimitação}\label{sec:regraEstrategia}

O receptor converte os cinco símbolos do gerador conforme a \autoref{tab:mapeamentoSimbolos}. Símbolos não mapeados são descartados apenas no caminho para o FPGA; a tabela é uma convenção do protótipo, não uma codificação de mercado.

\begin{table}[htb]
\caption{Mapeamento de símbolos do fluxo sintético}
\label{tab:mapeamentoSimbolos}
\centering
\begin{tabular}{lc@{\hspace{2cm}}lc}
\toprule
\textbf{Símbolo} & \textbf{Identificador} & \textbf{Símbolo} & \textbf{Identificador} \\ \midrule
AAPL & 0 & MSFT & 1 \\
NVDA & 2 & GOOGL & 3 \\
TSLA & 4 & & \\ \bottomrule
\end{tabular}
\fonte{Elaborado pelos autores conforme \codigo{fast_receiver.cpp}}
\end{table}

A regra retorna NOOP se um lado estiver vazio, se a melhor venda não superar a melhor compra ou se o \textit{spread} exceder 2,5000 unidades monetárias. Nos demais casos, retorna BUY para desbalanceamento maior ou igual a 500, SELL para valor menor ou igual a $-500$ e NOOP no intervalo restante. O \autoref{codigo:pseudocodigoEstrategia} explicita a ordem dessas condições e os valores de retorno implementados em \codigo{strategy.m}.

\begin{sourcecode}[htb]
\caption{Pseudocódigo da regra de decisão}
\label{codigo:pseudocodigoEstrategia}
\begin{lstlisting}[frame=single]
FUNCAO decidir(melhor_compra, qtd_compra,
               melhor_venda, qtd_venda,
               spread_1e4, desbalanceamento)
  acao <- NOOP

  SE qtd_compra = 0 OU qtd_venda = 0 ENTAO
    RETORNE acao
  FIM SE

  SE melhor_venda <= melhor_compra ENTAO
    RETORNE acao
  FIM SE

  SE spread_1e4 > 25000 ENTAO
    RETORNE acao
  FIM SE

  SE desbalanceamento >= 500 ENTAO
    acao <- BUY
  SENAO SE desbalanceamento <= -500 ENTAO
    acao <- SELL
  FIM SE

  RETORNE acao
FIM FUNCAO
\end{lstlisting}
\fonte{Elaborado pelos autores conforme \codigo{matlab/strategy.m}}
\end{sourcecode}

Os limiares foram escolhidos pelos autores para produzir ramos verificáveis com o fluxo sintético; não foram calibrados com dados históricos e não devem ser interpretados como estratégia financeira.

O receptor possui dois modos de execução. Sem a variável \codigo{HFT_FPGA_MMIO_BASE}, ele opera somente em software: conecta-se ao gerador, decodifica e imprime as mensagens. Com a variável definida, também abre a região MMIO, confere assinatura e geometria e publica eventos na fila TX. A separação entre os modos ajuda a distinguir problemas de protocolo de problemas da integração física.

O escopo termina na impressão da resposta. Não há canal de envio de ordens, conexão com bolsa, dados históricos, controle de posição, gestão de risco, multicast, carimbo de tempo de rede ou avaliação de resultado financeiro. Desse modo, a execução funcional valida a cadeia implementada, não uma operação de HFT completa.
